\documentclass[10pt,a4paper]{amsart}
\usepackage[margin=19mm]{geometry}
\usepackage{amsmath,amssymb,mathtools}
\usepackage[scr=rsfs,cal=euler]{mathalfa}
\usepackage{xfrac}
\usepackage[unicode,hidelinks,bookmarks=false]{hyperref}
\newcommand{\ie}{i.e.\ }
\newcommand{\cf}{cf.\ }
\newcommand{\resp}{respectively\ }
\newcommand{\Subsection}{\subsection}
\providecommand{\nobreakdash}{\nobreak\hbox{-}}
\setlength{\emergencystretch}{2em}
\multlinegap0pt
\usepackage{bm}
\usepackage{bbm}
\usepackage{dsfont}
\usepackage{xspace}
\usepackage{cancel}
\usepackage{mathtools}
\usepackage{amsthm}
\makeatletter
\@ifundefined{thm}{\newtheorem{thm}{Thm}}{}
\@ifundefined{lem}{\newtheorem{lem}[thm]{Lemma}}{}
\makeatother
\renewcommand{\le}{\leqslant}
\title[Dimension-dependent bounds for the SDIEP via phase optimisat]{Dimension-dependent bounds for the SDIEP via phase optimisation and Paley-type constructions}
\author{Tomasz Kania}
\date{Source: arXiv:2509.24079v1 (CC BY 4.0)}
\begin{document}
\maketitle

\noindent\emph{Source and licence.} Dimension-dependent bounds for the SDIEP via phase optimisation and Paley-type constructions. Version arXiv:2509.24079v1 (CC BY 4.0), distributed under the Creative Commons Attribution 4.0 International licence, as recorded in the arXiv licence metadata for this exact version and shown on the abstract page https://arxiv.org/abs/2509.24079v1. Full rights notice: licenses/arxiv-2509.24079-CC-BY-4.0.txt.\par\smallskip

\noindent\textit{Editorial scope.} Counted unit: the Lem whose source label is \texttt{lem:twophase} in the article (source0.tex lines 295--316, source label \texttt{lem:twophase}), delivered together with the article's own complete original proof of it (source0.tex lines 318--409, one \texttt{proof} environment of 92 lines). No mathematics is written, repaired or re-derived: statement and proof are the article's own text line for line. Required context retained uncounted: none. Every definition, display and label of those lines is retained unchanged. Omitted from this package and not redistributed: 1--294, 317--317, 410--756. Nothing inside a retained excerpt is omitted or altered: every retained body is delivered exactly as the article's lines are written, so no omission receipt of any kind accompanies this package. Citations: the retained excerpts carry no citation command at all, so no bibliography entry of the article is transcribed and no reference is redistributed. Reference adaptation: none was needed and none was made; every reference inside the delivered unit resolves inside it, so no unresolved reference or citation is produced. Printed slips kept literal: no printed word, symbol, hypothesis or number of the counted statement or of its proof was altered. Layout and class: the article's own class is \texttt{amsart} with packages including inputenc, fontenc, a4wide, scrextend, amssymb, amsthm, latexsym, mathrsfs, dsfont, euscript, esint, newtxtext, newtxmath, xcolor; the wrapper is the lane's pinned \texttt{amsart} editorial wrapper at 10pt on A4, which restates the theorem-like environments the retained text needs and adds the article's own macro definitions that the retained text needs (\texttt{\textbackslash le}), with extra wrapper packages (bm, bbm, dsfont, xspace, cancel, mathtools). The delivered numbering is the unit's own. Assets: no retained span contains a figure environment, a graphics-inclusion command, a PDF-page inclusion, a table or a TikZ picture; no image file is copied into this package, the package declares no asset dependency and no image file is redistributed with it. Licence: version arXiv:2509.24079v1 of this article is distributed under the Creative Commons Attribution 4.0 International licence, as recorded in the arXiv licence metadata for this exact version and shown on the abstract page https://arxiv.org/abs/2509.24079v1. A full rights notice is delivered at licenses/arxiv-2509.24079-CC-BY-4.0.txt. Characters: every retained excerpt is pure ASCII, so the delivered engine, XeLaTeX with its default Latin Modern text font, reports no missing character.
\begin{lem}[Two-phase covering radius on a grid]\label{lem:twophase}
Let $n':=n/\gcd(j,n)$ and set $S:=2\pi/n'$ and $L:=\pi/2$. All angles are considered reduced modulo $L$. Define
\[
\Delta^\star(n')\ :=\
\begin{cases}
\pi/n',     & n'\equiv 0\ (\mathrm{mod}\ 4),\\[1mm]
\pi/(2n'),  & n'\equiv 2\ (\mathrm{mod}\ 4),\\[1mm]
\pi/(4n'),  & n'\ \mathrm{odd}.
\end{cases}
\]
For $w^{(\phi)}_j(k)=\sqrt{\tfrac{2}{n}}\sin(\tfrac{2\pi jk}{n}+\phi)$ and
$\tilde w^{(\phi)}_j(k)=\sqrt{\tfrac{2}{n}}\cos(\tfrac{2\pi jk}{n}+\phi)$ one has
\[
\min_{\phi}\ \max\!\Big\{\,\|w^{(\phi)}_j\|_\infty,\ \|\tilde w^{(\phi)}_j\|_\infty\,\Big\}
\ =\ \sqrt{\tfrac2n}\ \cos\big(\Delta^\star(n')\big),
\]
and consequently
\[
\min_{\phi}\ \sup_{k,l}\ n\,\max\{|w^{(\phi)}_j(k)w^{(\phi)}_j(l)|,\ |\tilde w^{(\phi)}_j(k)\tilde w^{(\phi)}_j(l)|\}
\ =\ 2\cos^2\big(\Delta^\star(n')\big).
\]
\end{lem}

\begin{proof}
Let $A:=\{mL:\ m\in\mathbb Z\}$ with $L=\pi/2$. For any $x\in\mathbb R$,
\begin{equation}\label{eq:max-sin-cos}
\max\{|\sin x|,|\cos x|\}=\cos\big(\operatorname{dist}(x,A)\big).
\end{equation}
Indeed, reduce $x$ modulo $L$ to $y\in[0,L]$. Then
\[
|\cos x|=|\cos y|,\qquad |\sin x|=|\sin y|=|\cos(L-y)|.
\]
Hence $\max\{|\sin x|,|\cos x|\}=\max\{|\cos y|,|\cos(L-y)|\}=\cos\big(\min\{y,L-y\}\big)$.
However $\min\{y,L-y\}=\operatorname{dist}(x,A)$ by construction, giving \eqref{eq:max-sin-cos}.

\smallskip

Fix $j$ and write $n':=n/\gcd(j,n)$. The set of phases
\[
\Theta_j:=\Big\{\tfrac{2\pi jk}{n}\,:\ k=0,1,\dots,n-1\Big\}
\]
is a union of $\gcd(j,n)$ translates of the uniform grid $\{2\pi m/n':\ m=0,\dots,n'-1\}$. Since a~translation does not affect sup norms, it suffices to study the latter uniform grid. After adding a~free phase $\phi$ and reducing modulo $L$ we obtain
\[
X_\phi:=\Big\{\tfrac{2\pi m}{n'}+\phi\ \bmod L:\ m=0,\dots,n'-1\Big\}\ \subset \ \mathbb R/L\mathbb Z.
\]
This set is a coset (translate) of the cyclic subgroup $G:=\langle \overline{S}\rangle$ of the circle $\mathbb R/L\mathbb Z$, generated by the step $\overline{S}$ which is the class of $S:=2\pi/n'$ modulo $L$.

\smallskip

The order $r:=|G|$ is the least positive integer such that $r\overline{S}=\overline{0}$, i.e.
\[
rS\in L\mathbb Z \quad\Longleftrightarrow\quad \frac{2\pi r}{n'}=\frac{\pi}{2}\,q \ \text{ for some } q\in\mathbb Z
\quad\Longleftrightarrow\quad \frac{4r}{n'}=q\in\mathbb Z.
\]
Thus $r$ is the least positive integer with $n' \mid 4r$, which is
\[
r=\frac{n'}{\gcd(n',4)}.
\]
Consequently, the $r$ points of $G$ are equally spaced on the circle $\mathbb R/L\mathbb Z$ with spacing
\begin{equation}\label{eq:seff}
s_{\mathrm{eff}}=\frac{L}{r}=\frac{\pi}{2}\cdot\frac{\gcd(n',4)}{n'}
=
\begin{cases}
2\pi/n',   & n'\equiv 0\ (\mathrm{mod}\ 4),\\[1mm]
\pi/n',    & n'\equiv 2\ (\mathrm{mod}\ 4),\\[1mm]
\pi/(2n'), & n'\ \text{odd}.
\end{cases}
\end{equation}

\smallskip

For a fixed translate $X_\phi=\overline{\phi}+G$, the maximal distance from the origin $\overline{0}$ to $X_\phi$ is
\[
\rho(\phi):=\max_{x\in X_\phi}\operatorname{dist}(x,\overline{0}).
\]
Because $X_\phi$ is equally spaced with gap $s_{\mathrm{eff}}$, the circle $\mathbb R/L\mathbb Z$ is partitioned into $r$ closed arcs of length $s_{\mathrm{eff}}$ whose endpoints are the points of $X_\phi$. The point $\overline{0}$ lies in a unique such arc; its distance to the nearer endpoint is at most $s_{\mathrm{eff}}/2$, and equals $s_{\mathrm{eff}}/2$ precisely when $\overline{0}$ is at the midpoint of that arc. Therefore
\begin{equation}\label{eq:radius}
\min_{\phi}\ \rho(\phi)\ =\ \frac{s_{\mathrm{eff}}}{2}.
\end{equation}
Moreover, such a minimising shift $\phi$ exists: take any $x_0\in G$ and choose $\phi$ so that $\overline{\phi}+x_0=\overline{s_{\mathrm{eff}}/2}$.

\smallskip

For any $\phi$ we have, by \eqref{eq:max-sin-cos},
\[
\max\Big\{\|w^{(\phi)}_j\|_\infty,\ \|\tilde w^{(\phi)}_j\|_\infty\Big\}
=\sqrt{\tfrac{2}{n}}\ \max_{x\in X_\phi}\max\{|\sin x|,|\cos x|\}
=\sqrt{\tfrac{2}{n}}\ \max_{x\in X_\phi}\cos\big(\operatorname{dist}(x,\overline{0})\big).
\]
Since $\cos$ is decreasing on $[0,\pi/2]$, the inner maximum equals $\cos(d(\phi))$ where
\[
d(\phi):=\min_{x\in X_\phi}\operatorname{dist}(x,\overline{0}).
\]
Thus
\[
\min_{\phi}\ \max\Big\{\|w^{(\phi)}_j\|_\infty,\ \|\tilde w^{(\phi)}_j\|_\infty\Big\}
=\sqrt{\tfrac{2}{n}}\ \min_{\phi}\cos\big(d(\phi)\big)
=\sqrt{\tfrac{2}{n}}\ \cos\!\Big(\max_{\phi} d(\phi)\Big).
\]

Because $X_\phi$ is a translate of a cyclic subgroup $G$ of order $r=\frac{n'}{\gcd(n',4)}$ consisting of $r$ equally spaced points with gap $s_{\mathrm{eff}}=L/r$, the circle $\mathbb R/L\mathbb Z$ is partitioned into $r$ arcs of length $s_{\mathrm{eff}}$ with endpoints $X_\phi$. Placing $\overline{0}$ at the midpoint of one such arc yields 
\[
\max_\phi d(\phi)=\frac{s_{\mathrm{eff}}}{2}.
\]
Combining with \eqref{eq:seff} gives the stated $\sqrt{2/n}\,\cos(\Delta^\star(n'))$.

Finally, for either column in the pair, we obtain
\[
\min_{\phi}\ \sup_{k,l}\ n\,|q(k)q(l)|
\ \le\ \min_{\phi}\ n\,\|q\|_\infty^2
\ \le\ n\cdot \Big(\sqrt{\tfrac{2}{n}}\cos(\Delta^\star(n'))\Big)^2
=2\cos^2(\Delta^\star(n')).
\]
This upper bound is tight because, at the minimising shift, there are $k$ attaining the distance $\rho(\phi)=s_{\mathrm{eff}}/2$, hence $|q(k)|=\sqrt{2/n}\,\cos(\Delta^\star(n'))$ for one of the two vectors. Taking the maximum over the two vectors then attains $2\cos^2(\Delta^\star(n'))$ on some pair $(k,l)$, proving the second claim.
\end{proof}

\end{document}
